% Part: first-order-logic
% Chapter: beyond
% Section: intuitionistic-logic

\documentclass[../../../include/open-logic-section]{subfiles}

\begin{document}

\olfileid[hi]{fol}{byd}{il}

\olsection{अंतःप्रज्ञावादी तर्कशास्त्र}

द्वितीय-कोटि और उच्चतर-कोटि तर्कशास्त्र के विपरीत अंतःप्रज्ञावादी
प्रथम-कोटि तर्कशास्त्र चिरसम्मत रूप का एक प्रतिबंध है, जिसका उद्देश्य अधिक
``रचनात्मक'' प्रकार के विचार-विन्यास का प्रतिरूपण करना है। निम्न उदाहरण
उसकी कुछ मूल प्रेरणाएँ स्पष्ट कर सकते हैं।

मान लीजिए कि एक दिन कोई आपके पास आकर घोषणा करता है कि उसने एक प्राकृतिक
संख्या~$x$ निर्धारित की है जिसका गुण यह है: यदि $x$ अभाज्य है तो रीमान
परिकल्पना सत्य है, और यदि $x$ संयुक्त है तो रीमान परिकल्पना असत्य है। बहुत
अच्छी खबर!{} रीमान परिकल्पना सत्य है या नहीं, यह गणित के बड़े अनसुलझे
प्रश्नों में से एक है; और यहाँ उस व्यक्ति ने समस्या को मानो परिकलन की समस्या
में, अर्थात् किसी निश्चित संख्या के अभाज्य होने या न होने के निर्धारण में,
घटा दिया है।

$x$ का वह जादुई मान क्या है? वह इसे इस प्रकार बताता है: $x$ वह प्राकृतिक
संख्या है जो रीमान परिकल्पना सत्य होने पर $7$, और अन्यथा $9$ के बराबर है।

क्रोधित होकर आप अपना पैसा वापस माँगते हैं। चिरसम्मत दृष्टि से ऊपर का वर्णन
वास्तव में $x$ का एक अद्वितीय मान निर्धारित करता है; किंतु आप वास्तव में
$x$ का ऐसा मान चाहते हैं जो \emph{स्पष्टतः} दिया गया हो।

एक दूसरा, शायद कम कृत्रिम उदाहरण लें। हम जानते हैं कि किसी अपरिमेय संख्या
को परिमेय घात तक उठाकर परिमेय परिणाम पाया जा सकता है; जैसे,
$\sqrt{2}^2 = 2$। यह कम स्पष्ट है कि किसी अपरिमेय संख्या को
\emph{अपरिमेय} घात तक उठाकर परिमेय परिणाम पाया जा सकता है या नहीं। निम्न
प्रमेय इसका सकारात्मक उत्तर देता है:

\begin{thm}
ऐसी अपरिमेय संख्याएँ $a$ और $b$ हैं जिनके लिए $a^b$ परिमेय है।
\end{thm}

\begin{proof}
$\sqrt{2}^{\sqrt{2}}$ पर विचार कीजिए। यदि यह परिमेय है, तो काम पूरा है:
हम $a = b = \sqrt{2}$ ले सकते हैं। अन्यथा यह अपरिमेय है। तब
\[
(\sqrt{2}^{\sqrt{2}})^{\sqrt{2}} = \sqrt{2}^{\sqrt{2} \cdot
  \sqrt{2}} = \sqrt{2}^2 = 2,
\]
जो निश्चय ही परिमेय है। अतः इस स्थिति में $a$ को
$\sqrt{2}^{\sqrt{2}}$ और $b$ को~$\sqrt 2$ लें।
\end{proof}

क्या यह वैध प्रमाण है? अधिकांश गणितज्ञ मानते हैं कि है। फिर भी इसमें कुछ
असंतोषजनक है: हमने एक निश्चित गुण वाले वास्तविक संख्याओं के युग्म का
अस्तित्व सिद्ध किया, पर यह नहीं बता सके कि वह \emph{कौन-सा} युग्म है। उसी
परिणाम को इस प्रकार भी सिद्ध किया जा सकता है कि युग्म $a$, $b$ प्रमाण में
\emph{दिया} हो: $a = \sqrt{3}$ और $b = \log_3 4$ लें। तब
\[
a^b = \sqrt{3}^{\log_3 4} = 3^{1/2 \cdot \log_3 4} = (3^{\log_3
  4})^{1/2} = 4^{1/2}= 2,
\]
क्योंकि $3^{\log_3 x} = x$।

अंतःप्रज्ञावादी तर्कशास्त्र ऐसे विचार-विन्यास का प्रतिरूपण करने के लिए बना
है जिसमें पहले प्रमाण जैसी चाल मान्य नहीं है। किसी $x$ का, जो~$!A(x)$ को
संतुष्ट करता हो, अस्तित्व सिद्ध करने का अर्थ है कि आपको कोई निश्चित~$x$ और उसके
$!A$ को संतुष्ट करने का प्रमाण देना होगा, जैसा दूसरे प्रमाण में है। $!A$ या
$!B$ सत्य है, यह सिद्ध करने के लिए इनमें से किसी एक को सिद्ध कर सकना आवश्यक
है।

औपचारिक रूप से, प्रथम-कोटि तर्कशास्त्र के किसी !!a{derivation}-तंत्र को एक
निश्चित ढंग से प्रतिबंधित करने पर अंतःप्रज्ञावादी प्रथम-कोटि तर्कशास्त्र
मिलता है। इसी प्रकार द्वितीय-कोटि और उच्चतर-कोटि तर्कशास्त्र के भी
अंतःप्रज्ञावादी रूप हैं। गणितीय दृष्टि से ये केवल औपचारिक निगमन-तंत्र हैं,
पर जैसा कहा गया, इनका उद्देश्य एक प्रकार के गणितीय विचार-विन्यास का
प्रतिरूपण करना है। इसे वह विचार-विन्यास माना जा सकता है जिसे गणित का कोई
निश्चित दार्शनिक मत (जैसे ब्राउअर का अंतःप्रज्ञावाद) उचित ठहराता है; इसे
अधिक ``मूर्त'' और संतोषजनक गणितीय विचार-विन्यास (बिशप के रचनावाद की दिशा
में) माना जा सकता है; और इस पर विवाद किया जा सकता है कि औपचारिक वर्णन
अनौपचारिक प्रेरणा को ग्रहण करता है या नहीं। पर हमारे दार्शनिक मत कुछ भी हों,
हम अंतःप्रज्ञावादी तर्कशास्त्र का अध्ययन औपचारिक रूप से प्रस्तुत तर्कशास्त्र
के रूप में कर सकते हैं; और कारण जो भी हों, अनेक गणितीय तर्कशास्त्री ऐसा
करना रोचक पाते हैं।

अंतःप्रज्ञावादी संयोजकों की एक अनौपचारिक रचनात्मक व्याख्या है, जिसे सामान्यतः
BHK व्याख्या कहा जाता है (ब्राउअर, हेटिंग और कोल्मोगोरोव के नाम पर)। वह इस
प्रकार है: $!A
\land !B$ का प्रमाण, $!A$ के प्रमाण के साथ युग्मित $!B$ के
प्रमाण से बनता है; $!A \lor !B$ का प्रमाण या तो $!A$ का प्रमाण है या $!B$
का, और हमें स्पष्ट सूचना होती है कि इनमें से कौन-सी स्थिति है; $!A \lif !B$
का प्रमाण ऐसी प्रक्रिया है जो $!A$ के प्रमाण को~$!B$ के प्रमाण में बदलती
है; $\lforall[x][!A(x)]$ का प्रमाण ऐसी प्रक्रिया है जो $!A(x)$ का प्रमाण
~$x$ के प्रत्येक मान के लिए लौटाती है; और $\lexists[x][!A(x)]$ का प्रमाण~$x$
के किसी मान और उस मान द्वारा~$!A$ को संतुष्ट करने के प्रमाण से बनता है। इस
व्याख्या को ``करी--हावर्ड समरूपता'' अथवा ``!!{formula}s-को-प्रकार मानने का
प्रतिमान'' कहलाने वाले संगणनात्मक पदों में भी कह सकते हैं: किसी !!a{formula}
को एक निश्चित प्रकार के आँकड़ा-प्रकार का विनिर्देशन और प्रमाणों को उन
आँकड़ा-प्रकारों की संगणनात्मक वस्तुएँ मानिए, जिनसे हम देख सकते हैं कि संगत
!!{formula} सत्य है।

अंतःप्रज्ञावादी तर्कशास्त्र को प्रायः मध्यनिषेध नियम ``घटाकर'' बचा हुआ
चिरसम्मत तर्कशास्त्र माना जाता है। निम्न प्रमेय इस कथन को अधिक सटीक बनाता
है।
\begin{thm}
अंतःप्रज्ञावादी दृष्टि से निम्न अभिगृहीत-प्रतिरूप समतुल्य हैं:
\begin{enumerate}
\item $(\lnot !A \lif \lfalse) \lif !A$.
\item $!A \lor \lnot !A$
\item $\lnot \lnot !A \lif !A$
\end{enumerate}
\end{thm}
अन्य दोनों में से किसी एक द्वारा एक प्रतिरूप की निष्पत्तियाँ प्राप्त करना
अंतःप्रज्ञावादी तर्कशास्त्र का अच्छा अभ्यास है।

ब्राउअर के अंतःप्रज्ञावाद के औपचारीकरण के रूप में प्रस्तुत
अंतःप्रज्ञावादी प्रतिज्ञप्ति-तर्कशास्त्र के प्रथम निगमन-तंत्र कोल्मोगोरोव,
ग्लिवेंको और हेटिंग ने स्वतंत्र रूप से दिए। अंतःप्रज्ञावादी प्रथम-कोटि
तर्कशास्त्र (और अंतःप्रज्ञावादी गणित के कुछ भागों) का प्रथम औपचारीकरण हेटिंग
ने किया। यद्यपि चिरसम्मत रूप से वैध अनेक प्रतिरूप अंतःप्रज्ञावादी रूप से
वैध नहीं हैं, अनेक वैध हैं।

\emph{द्विनिषेध अनुवाद} चिरसम्मत और अंतःप्रज्ञावादी तर्कशास्त्र के बीच एक
महत्त्वपूर्ण संबंध बताता है। इसे निम्न प्रकार आगमनिक रूप से परिभाषित किया
जाता है ($!A^N$ को चिरसम्मत !!{formula}~$!A$ का ``अंतःप्रज्ञावादी'' अनुवाद
समझिए):
\begin{align*}
!A^N & \ident \lnot\lnot !A \quad \text{परमाण्विक !!{formula}s $!A$ के लिए} \\
(!A \land !B)^N & \ident (!A^N \land !B^N) \\
(!A \lor !B)^N & \ident  \lnot\lnot (!A^N \lor !B^N) \\
(!A \lif !B)^N & \ident (!A^N \lif !B^N) \\
(\lforall[x][!A])^N & \ident \lforall[x][!A^N] \\
(\lexists[x][!A])^N & \ident \lnot\lnot\lexists[x][!A^N]
\end{align*}
कोल्मोगोरोव और ग्लिवेंको ने प्रतिज्ञप्ति-तर्कशास्त्र के लिए इस अनुवाद के रूप
दिए; विधेय-तर्कशास्त्र के लिए गोडेल (G\"odel) और जेंट्सेन ने इसे स्वतंत्र रूप से दिया।
हमारे पास निम्न प्रमेय है:

\begin{thm}
\begin{enumerate}
\item $!A \liff !A^N$ चिरसम्मत रूप से प्रमाण्य है।
\item यदि $!A$ चिरसम्मत रूप से प्रमाण्य है, तो $!A^N$ अंतःप्रज्ञावादी रूप से
  प्रमाण्य है।
\end{enumerate}
\end{thm}

अब निम्न संवाद की कल्पना कीजिए। चिरसम्मत गणितज्ञ: ``मैंने $!A$ सिद्ध कर
दिया!'' अंतःप्रज्ञावादी गणितज्ञ: ``तुम्हारा प्रमाण वैध नहीं है। वास्तव में
तुमने $!A^N$ सिद्ध किया है।'' चिरसम्मत गणितज्ञ: ``मुझे मंज़ूर है!''
चिरसम्मत गणितज्ञ की दृष्टि में अंतःप्रज्ञावादी केवल बाल की खाल निकाल रहा है,
क्योंकि दोनों समतुल्य हैं। किंतु अंतःप्रज्ञावादी आग्रह करता है कि उनमें अंतर
है।

ध्यान दें कि ऊपर का अनुवाद केवल शुद्ध तर्कशास्त्र से संबंधित है; वह इस
प्रश्न का उत्तर नहीं देता कि चिरसम्मत और अंतःप्रज्ञावादी गणित के लिए उपयुक्त
\emph{अतार्किक} अभिगृहीत कौन-से हैं, अथवा उनके बीच क्या संबंध है। किंतु ऊपर
के प्रमेय का निम्न छोटा-सा विस्तार कुछ उपयोगी सूचना देता है:

\begin{thm}
यदि $\Gamma$, $!A$ को चिरसम्मत रूप से सिद्ध करता है, तो $\Gamma^N$, $!A^N$
को अंतःप्रज्ञावादी रूप से सिद्ध करता है।
\end{thm}

दूसरे शब्दों में, यदि कुछ परिकल्पनाओं से $!A$ चिरसम्मत रूप से प्रमाण्य है,
तो उनके द्विनिषेध अनुवादों से $!A^N$ प्रमाण्य है।

यह दिखाने के लिए कि कोई वाक्य अथवा प्रतिज्ञप्तिक !!{formula}
अंतःप्रज्ञावादी रूप से वैध है, केवल एक प्रमाण देना होता है। किंतु यह कैसे
दिखाएँ कि वह वैध नहीं है? इसके लिए हमें ऐसा अर्थविज्ञान चाहिए जो सुदृढ़ और
अधिमानतः पूर्ण हो। क्रिप्के द्वारा दिया अर्थविज्ञान इस काम के लिए उपयुक्त
है।

हम चिरसम्मत तर्कशास्त्र वाला वही खेल खेल सकते हैं: अर्थविज्ञान परिभाषित करें
और सुदृढ़ता तथा पूर्णता सिद्ध करें। फिर भी निम्न भेद ध्यान देने योग्य है।
चिरसम्मत तर्कशास्त्र में अर्थविज्ञान एक अर्थ में ``स्पष्ट'' था और संयोजकों
के अर्थ में अंतर्निहित था। क्रिप्के अर्थविज्ञान के लिए कुछ सहज प्रेरणा दी
जा सकती है, पर उसमें वैसी ही अनिवार्यता की अनुभूति नहीं होती। इसके अतिरिक्त
चिरसम्मत !!{structure} की धारणा एक स्वाभाविक गणितीय धारणा है; इसलिए हम
!!a{structure} की धारणा को चिरसम्मत प्रथम-कोटि तर्कशास्त्र के अध्ययन का
साधन मान सकते हैं, अथवा यह मान सकते हैं कि चिरसम्मत प्रथम-कोटि तर्कशास्त्र
से गणितीय !!{structure}s समझ में आती हैं। इसके विपरीत क्रिप्के
!!{structure}s केवल तार्किक रचनाएँ मानी जा सकती हैं; उनमें स्वतंत्र गणितीय रुचि दिखाई नहीं
देती।

किसी प्रतिज्ञप्तिक भाषा के लिए क्रिप्के !!{structure}~$\mModel M = \tuple{W, R, V}$
में एक समुच्चय~$W$, आंशिक क्रम~$R$ जो~$W$ पर है और जिसका न्यूनतम !!{element}
है, तथा~$W$ के !!{element}s और प्रतिज्ञप्तिक चरों के बीच एक ``क्रम-संरक्षी''
नियतन होता है। सहज धारणा यह है कि $W$ के !!{element}s ``जगतों''
अथवा ``ज्ञान-अवस्थाओं'' को निरूपित करते हैं; अवयव $v \geq
u$,~$u$ की एक
``संभावित भावी अवस्था'' को निरूपित करता है; और~$u$ को नियत प्रतिज्ञप्तिक चर
वे प्रतिज्ञप्तियाँ हैं जिन्हें अवस्था~$u$ में सत्य ज्ञात किया जाता है। तब
बलन-संबंध $\mSat{M}{!A}[w]$ इस संबंध को भाषा के मनचाहे !!{formula}s तक
विस्तृत करता है; $\mSat{M}{!A}[w]$ को ``$!A$ अवस्था~$w$ में सत्य है'' पढ़ें।
यह संबंध निम्न प्रकार आगमनिक रूप से परिभाषित है:
\begin{enumerate}
\item $\mSat{M}{\Obj p_i}[w]$ तभी और केवल तभी जब $\Obj p_i$,~$w$ को नियत
  प्रतिज्ञप्तिक चरों में से एक हो।
\item $\mSat/{M}{\lfalse}[w]$।
\item $\mSat{M}{(!A \land !B)}[w]$ तभी और केवल तभी जब $\mSat{M}{!A}[w]$ और
  $\mSat{M}{!B}[w]$।
\item $\mSat{M}{(!A \lor !B)}[w]$ तभी और केवल तभी जब $\mSat{M}{!A}[w]$ या
  $\mSat{M}{!B}[w]$।
\item $\mSat{M}{(!A \lif !B)}[w]$ तभी और केवल तभी जब $w' \geq w$ और
  $\mSat{M}{!A}[w']$ होने पर सदा $\mSat{M}{!B}[w']$ हो।
\end{enumerate}
प्रतिउदाहरण देने वाली क्रिप्के !!{structure} बनाकर यह दिखाने का प्रयास एक
अच्छा अभ्यास है कि $\lnot (p \land q) \lif
(\lnot p \lor \lnot q)$
अंतःप्रज्ञावादी रूप से वैध नहीं है।

\end{document}
